% =============================================================================
%  resumen.tex
% =============================================================================

\chapter*{Resumen}
\addcontentsline{toc}{chapter}{Resumen}
\markboth{Resumen}{}

Los sistemas de clasificación automática de texto basados en procesamiento de
lenguaje natural se emplean cada vez con mayor frecuencia en procesos de alto
impacto social, como el filtrado de currículums, la moderación de contenido o
la categorización de solicitudes en plataformas digitales. Aunque suelen
presentarse como herramientas neutrales, distintas investigaciones han mostrado
que pueden reproducir o amplificar desigualdades sociales preexistentes. Este
trabajo aborda, desde un enfoque sociotécnico, la pregunta de en qué medida
estos sistemas presentan sesgos asociados al \emph{registro lingüístico} o
\emph{sociolecto}, esto es, a variaciones de estilo que no afectan el contenido
sustantivo de la información evaluada.

La metodología combina revisión documental, recolección de percepciones
mediante encuesta a personas jóvenes de distintos contextos socioeconómicos en
México y experimentación computacional sobre uno o más modelos de clasificación
textual. A partir de los patrones identificados se construye un conjunto de
textos controlados que mantienen constante el contenido informativo y varían
únicamente el registro lingüístico, los cuales se utilizan como
\emph{benchmark}. Sobre los resultados experimentales se define y calcula un
indicador preliminar, el \emph{Índice de Neutralidad Sociolectal} (\INS),
orientado a medir el grado de consistencia con que un sistema responde a estas
variaciones.

\paragraph{Palabras clave.} sesgo algorítmico, procesamiento de lenguaje
natural, discriminación lingüística, sociolectos, evaluación sociotécnica,
equidad algorítmica.

\vspace{1.5cm}

\section*{Abstract}
\addcontentsline{toc}{section}{Abstract}

Automatic text classification systems based on natural language processing are
increasingly deployed in socially consequential pipelines, such as résumé
screening, content moderation, and request triage in digital platforms.
Although often framed as neutral tools, prior work has shown that they may
reproduce or amplify pre-existing social inequalities. This thesis addresses,
from a sociotechnical perspective, the extent to which such systems exhibit
biases associated with linguistic \emph{register} or \emph{sociolect}, that
is, stylistic variation that does not affect the substantive content of the
information being evaluated.

The methodology combines a documentary review, a perception survey administered
to young Mexicans from different socioeconomic backgrounds, and a computational
experiment over one or more text classification models. From the identified
patterns, a controlled benchmark is built in which informational content is
held constant while linguistic register is varied. Over the experimental
results, a preliminary indicator is defined and computed, the
\emph{Sociolectal Neutrality Index} (\INS), aimed at measuring the consistency
of system responses across such variations.

\paragraph{Keywords.} algorithmic bias, natural language processing, linguistic
discrimination, sociolects, sociotechnical evaluation, algorithmic fairness.

\cleardoublepage
